\begin{center}
ABSTRACT
\end{center}

\bigskip

\begin{center}
SIRENS OF THE SWARM:\\
REVEALING BINARY BLACK HOLE MERGERS AND SUPERNOVA EXPLOSIONS IN ACTIVE GALACTIC NUCLEI WITH OBSERVATIONAL PREDICTIONS FOR THE AGE OF MULTI-MESSENGER ASTRONOMY \\
\bigskip
BY\\
\bigskip
HARRISON EDWARD COOK, B.A.
\end{center}

\bigskip

\begin{center}
Doctor of Philosophy\\
New Mexico State University\\
Las Cruces, New Mexico, 2022\\
Dr. Wladimir Lyra, Advisor
\end{center}

\bigskip
\bigskip




% ABSTRACT (350 WORDS)
Active galactic nuclei host a supermassive black hole, a swarm of stars and black holes in a nuclear star cluster, and most critically, an accretion disk.
Some of those stars and black holes will lie on orbits aligned with the disk, while others will have orbits that cross through it.
Gas drag from the disk produces torques on their orbits that will align more orbits with the disk and once their orbits are sufficiently circular will cause embedded objects to migrate.
Since more massive objects migrate faster, they can sweep up lighter black holes to form binaries on their way towards migration traps, where inwards and outwards migration meet, which can drastically increase binary formation rates.
Experiencing more gas effects and dynamical encounters with other objects can cause these binaries to merge quickly relative to systems without such mechanisms.
Once close enough, runaway orbital shrinkage will occur until they emit a burst of energy in the form of gravitational waves as they warp space-time in the final moments before coalescing into a single object that can be detected by ground-based gravitational wave observatories.
The environment imprints unique signatures on the underlying black hole population that can be used to probe the conditions of the disk and star cluster with enough events.
Black holes in this binary formation channel can participate in many merger events that may result in intermediate mass black holes due to their high concentration.
Since the disk aligns the orbits and spins of embedded black holes over time and increases their mass through hierarchical events, the channel naturally replicates an anti-correlation present in the observed gravitational wave events where mergers between unequal mass black holes tend to have their spins more aligned with the binary's orbital angular momentum.
In addition to the gravitational waves, the shocks from the remnant interacting with the gas disk may produce a bright electromagnetic counterpart.
The stars of the nuclear star cluster experience the same orbital changes caused by the disk, and massive stars may detonate in supernova explosions.
While these light curves are well-studied and documented, the dense disk causes these to differ from naked supernovae as suggested by studies of those that must expand into dense envelopes expelled by strong winds of their progenitor stars.
As such, these transients could become sources of contamination in searches for merger counterparts without understanding their characteristics.
Additionally, while the accretion disk is an important factor for these phenomena, the details of their structures and conditions are not well constrained.
Therefore, we can use these multi-messenger transient events to learn about the disk that produces them.
In this dissertation I will show work I have done to explore this dynamic environment.
I used Monte Carlo simulations to study populations of black hole mergers to understand how variations to initial conditions and physical processes affect their mass ratios and effective spins.
I will also show hydrodynamic simulations of embedded supernovae and their light curves produced in post-processing radiative transfer calculations.

